\begin{abstract}\noindent
We study a hybrid real-estate investment structure in which the rental cash
flows of a Paris residential building are partitioned into a senior /
mezzanine / equity tranche stack akin to a synthetic CDO, and compare its
risk-return profile to two classical benchmarks: outright single-apartment
ownership (Model A) and a fully pooled SAS share (Model B). All market
parameters are calibrated on French sources: \dataParisN{} quarterly observations of the
Notaires-INSEE price index (\dataParisStart{}--\dataParisEnd{}), the INSEE
\textit{Indice de Référence des Loyers}, and the French ten-year sovereign
yield (FRED \texttt{IRLTLT01FRM156N}). Default times are simulated under
three benchmark credit models — the Gaussian one-factor copula
\autocite{Li2000, Vasicek2002}, the Student-t copula
\autocite{EmbrechtsMcNeilStraumann2002} and a Cox doubly-stochastic
intensity \autocite{DuffieSingleton2003} — coupled with a Vasicek short rate
\autocite{Vasicek1977} and a geometric Brownian motion (with a Merton
jump-diffusion extension \autocite{Merton1976}) for the property value. The
tranche cash flows follow the Andersen-Sidenius-Basu (2003) waterfall and
are priced by Monte Carlo with antithetic variance reduction and a Sobol
QMC variant. On the Paris-intermediate scenario the senior fair value drops from
\seniorFairGauss{} of par under the Gaussian copula to
\seniorFairStudent{} under the Student-t copula, and the senior fair
coupon widens by roughly twelve basis points --- a numerical
illustration of the tail-dependence critique levelled at Li-style
models post-2008 \autocite{DonnellyEmbrechts2010, MacKenzieSpears2014}.
\end{abstract}
